\section{Introduction}
\label{sec:intro}

Suppose we hand a document to an AI and ask it to \emph{review} the writing, then hand
both the document and that review back and ask it to \emph{revise}, and repeat. Does the
process ever settle? This is not an idle question: automated editing assistants, agentic
writing pipelines, and self-improving document workflows all run some version of this loop,
often with an implicit ``keep going until it stabilizes'' stopping rule. Whether such a loop
\emph{terminates}---and whether its limit is a genuinely better document or a degenerate
attractor---determines whether it is safe to run to ``stability.''

\para{The loop is a dynamical system.} Following recent work that views recursive text
generation as a map over documents~\citep{wu2026converge}, we treat a document version
$V_t$ as a state and the two-stage operator as $V_{t+1} = F(V_t)$ with
$F = \Revise(\cdot, \Review(\cdot))$. A \emph{fixed point} is a version the loop no longer
changes ($V_{t+1} = V_t$), or, more realistically, a soft region where consecutive versions
differ by a tiny, stable edit distance. The natural hypothesis---and the folk intuition
behind ``revise until it stops changing''---is that the loop contracts toward such a fixed
point.

\para{The gap.} The closest prior work,~\citet{wu2026converge}, studies a
\emph{single-stage} self-refine operator (revise directly, with no separate critique) and
finds it relaxes cleanly to a soft fixed point. The canonical Self-Refine
loop~\citep{madaan2023selfrefine} \emph{fuses} critique and revision in one prompt. Neither
isolates the effect of a \emph{decoupled, freshly regenerated review at every step}. This
matters because a review stage is not neutral: a reviewer is instructed to \emph{find
weaknesses}, and a capable model can always find some. The skeptical
literature---self-bias amplification~\citep{xu2024pride}, feedback
friction~\citep{feedbackfriction2025}, and recursion-induced
degradation~\citep{shumailov2023curse}---predicts recursion need not be benign, but does not
measure the review$\to$revise fixed-point question directly.

\para{Our approach.} We run the two-stage loop for $T=8$ iterations on 12 fresh 2025-era
scientific abstracts, logging every version $V_t$ and every review $R_t$, and measure
convergence with the exact metrics of~\citet{wu2026converge}: normalized Levenshtein edit
distance $\dt$, exact and approximate fixed-point rates, exponential-relaxation fits, plus
semantic drift and an LLM-as-judge of quality and meaning preservation. Crucially, we run a
\emph{single-stage} self-refine baseline on the same documents as a positive control, and
ablate seven factors (temperature, cross-model reviewer, review framing, genre, and length
guidance) to find what governs convergence.

\para{What we find.} The two-stage loop does \emph{not} converge. Per-step change plateaus
at a large residual floor (mean final $\Delta \approx 0.38$, an order of magnitude above the
$0.02$ approximate-fixed-point threshold; $t(11)=11.3$, $p=2.2\times10^{-7}$), and $0\%$ of
trajectories reach a fixed point. The single-stage baseline, on the \emph{same} documents,
does converge (final $\Delta \approx 0.05$, $42\%$ approximate fixed points, fit
$R^2=0.98$)---so the harness detects convergence when it exists, and the explicit review
stage is the decisive difference (Cohen's $d_z \approx -3.3$, the largest effect in the
study). The mechanism is runaway growth: a never-satisfied reviewer emits fresh demands
every round, the reviser satisfies them by \emph{adding} text, documents inflate $\sim 7\times$
(a 1.5\,KB abstract becomes a 10\,KB mini-paper), drift from the original (cosine
$0.97\to0.77$), and lose fidelity ($\sim 4$ fabricated claims per document)---while
isolated-quality judge scores barely move ($7.5\to 8.0$), so the loop looks locally like
improvement while globally diverging.

\para{Contributions.}
\begin{itemize}[leftmargin=*,itemsep=0pt,topsep=0pt]
    \item We formalize and directly test the decoupled review$\to$revise loop as a
    dynamical system, and show it does \emph{not} reach a fixed point: mean final
    $\Delta \approx 0.38$ and $0\%$ convergence across all two-stage conditions
    (\secref{sec:results}).
    \item We identify the explicit review stage as the cause via a same-document single-stage
    positive control that \emph{does} converge ($42\%$ approximate fixed points,
    $R^2=0.98$), the dominant effect in a seven-way ablation ($d_z\approx-3.3$).
    \item We characterize the failure mechanism---never-satisfied-reviewer-driven length
    inflation ($\sim 7\times$) that trades meaning preservation ($10\to 6.2$) for flat,
    slightly-rising isolated quality scores---and show it is not a truncation artifact.
    \item We draw a practical implication: ``run until it stabilizes'' is unsafe for
    review$\to$revise loops; a fixed point requires either dropping the review stage or
    adding an explicit stop criterion.
\end{itemize}

The rest of the paper reviews related work (\secref{sec:related}), formalizes the operator
and setup (\secref{sec:method}), presents results and ablations (\secref{sec:results}),
discusses interpretation and limitations (\secref{sec:discussion}), and concludes
(\secref{sec:conclusion}).
